\documentclass[../../main.tex]{subfiles}
\begin{document}
\section{含时向量场}\label{section:ch09-含时向量场}

迄今为止我们遇到的所有微分方程组都是自治的，即当它们写成\eqref{align:ch09-01}的形式时，右端的函数 $V^i$ 不显式依赖自变量 $t$（见附录 D）。然而，非自治 ODE 确实在流形理论中出现，因此值得探讨本章的结果如何推广到此情形。我们将仅在第 22 章使用此理论。

设 $M$ 是光滑流形。$M$ 上的\textbf{含时向量场}是连续映射 $V\colon J\times M\to TM$，其中 $J\subseteq\mathbb{R}$ 是区间，使得对每个 $(t;p)\in J\times M$，$V(t;p)\in T_pM$。这意味着对每个 $t\in J$，由 $V_t(p)=V(t;p)$ 定义的映射 $V_t\colon M\to TM$ 是 $M$ 上的向量场。若 $V$ 是 $M$ 上的含时向量场，则 $V$ 的\textbf{积分曲线}是可微曲线 $\gamma\colon J_0\to M$，其中 $J_0$ 是包含于 $J$ 的区间，满足对所有 $t\in J_0$，
\begin{align*}
\gamma'(t)=V(t;\gamma(t)).
\end{align*}

每个普通向量场 $X\in\mathfrak{X}(M)$ 都通过令 $V(t;p)=X_p$ 确定一个定义在 $\mathbb{R}\times M$ 上的含时向量场。（考虑定义在 $\mathbb{R}\times M$ 的更一般开子集上的含时向量场偶尔是有用的；但为简单起见我们仅限于乘积集 $J\times M$，将更一般情形下结果需如何修改留给感兴趣的读者。）

含时向量场不一定生成流，因为从同一点但不同时刻出发的两条积分曲线可能沿不同路径，而流经过给定点的所有积分曲线有相同的像。作为流基本定理的替代，我们有如下定理。

\begin{theorem}\label{theorem:ch09-48}[含时流的基本定理]
设 $M$ 是光滑流形，$J\subseteq\mathbb{R}$ 是开区间，$V\colon J\times M\to TM$ 是 $M$ 上的光滑含时向量场。存在开子集 $E\subseteq J\times J\times M$ 和光滑映射 $\Psi\colon E\to M$，称为 $V$ 的\textbf{含时流}，具有如下性质：
\begin{enumerate}[(a)]
\item 对每个 $t_0\in J$ 和 $p\in M$，集合 $E(t_0;p)=\{t\in J:(t;t_0;p)\in E\}$ 是包含 $t_0$ 的开区间，由 $\Psi^{(t_0;p)}(t)=\Psi(t;t_0;p)$ 定义的光滑曲线 $\Psi^{(t_0;p)}\colon E(t_0;p)\to M$ 是具有初始条件 $\Psi^{(t_0;p)}(t_0)=p$ 的唯一极大积分曲线。
\item 若 $t_1\in E(t_0;p)$ 且 $q=\Psi^{(t_0;p)}(t_1)$，则 $E(t_1;q)=E(t_0;p)$ 且 $\Psi^{(t_1;q)}=\Psi^{(t_0;p)}$。
\item 对每个 $(t_1;t_0)\in J\times J$，集合 $M_{t_1;t_0}=\{p\in M:(t_1;t_0;p)\in E\}$ 在 $M$ 中开，由 $\Psi_{t_1;t_0}(p)=\Psi(t_1;t_0;p)$ 定义的映射 $\Psi_{t_1;t_0}\colon M_{t_1;t_0}\to M$ 是从 $M_{t_1;t_0}$ 到 $M_{t_0;t_1}$ 的微分同胚，其逆为 $\Psi_{t_0;t_1}$。
\item 若 $p\in M_{t_1;t_0}$ 且 $\Psi_{t_1;t_0}(p)\in M_{t_2;t_1}$，则 $p\in M_{t_2;t_0}$ 且
\begin{align}
\Psi_{t_2;t_1}\circ\Psi_{t_1;t_0}(p)=\Psi_{t_2;t_0}(p). \label{align:ch09-18}
\end{align}
\end{enumerate}
\end{theorem}

\begin{proof}
这可通过遵循\cref{theorem:ch09-12}证明的梗概来证明，用\cref{theorem:appD-06}代替\cref{theorem:appD-01}。然而，用如下技巧将其约化为不含时情形要快得多。

考虑 $J\times M$ 上由
\begin{align*}
\widetilde{V}(s;p)=\left(\frac{\partial}{\partial s}\bigg|_s,\;V(s;p)\right)
\end{align*}
定义的光滑向量场 $\widetilde{V}$，其中 $s$ 是 $J\subseteq\mathbb{R}$ 上的标准坐标，我们像通常一样将 $T_{(s;p)}(J\times M)$ 等同于 $T_sJ\oplus T_pM$（见\cref{proposition:ch03-14}）。设 $\widetilde{\theta}\colon\widetilde{\mathscr{D}}\to J\times M$ 表示 $\widetilde{V}$ 的流。若将 $\widetilde{\theta}$ 的分量函数写为
\begin{align*}
\widetilde{\theta}(t;(s;p))=\bigl(\alpha(t;(s;p)),\beta(t;(s;p))\bigr),
\end{align*}
则 $\alpha\colon\widetilde{\mathscr{D}}\to J$ 和 $\beta\colon\widetilde{\mathscr{D}}\to M$ 满足
\begin{align*}
\frac{\partial\alpha}{\partial t}(t;(s;p))=1,\quad \alpha(0;(s;p))=s,\\
\frac{\partial\beta}{\partial t}(t;(s;p))=V\bigl(\alpha(t;(s;p));\beta(t;(s;p))\bigr),\quad \beta(0;(s;p))=p.
\end{align*}
由此立刻得 $\alpha(t;(s;p))=t+s$，因而 $\beta$ 满足
\begin{align}
\frac{\partial\beta}{\partial t}(t;(s;p))=V\bigl(t+s;\beta(t;(s;p))\bigr). \label{align:ch09-19}
\end{align}

设 $E$ 为 $\mathbb{R}\times J\times M$ 中由
\begin{align*}
E=\{(t;t_0;p):(t-t_0;(t_0;p))\in\widetilde{\mathscr{D}}\}
\end{align*}
定义的子集。显然，由于 $\widetilde{\mathscr{D}}$ 开，$E$ 在 $\mathbb{R}\times J\times M$ 中开。此外，由于 $\alpha$ 将 $\widetilde{\mathscr{D}}$ 映入 $J$，若 $(t;t_0;p)\in E$，则 $t=\alpha(t-t_0;(t_0;p))\in J$，这推出 $E\subseteq J\times J\times M$。

每个集合 $M_{t_1;t_0}=\{p\in M:(t_1;t_0;p)\in E\}$ 在 $M$ 中开，这是 $E$ 开的直接推论。

现定义 $\Psi\colon E\to M$ 为
\begin{align*}
\Psi(t;t_0;p)=\beta(t-t_0;(t_0;p)).
\end{align*}
则 $\Psi$ 光滑（因为 $\beta$ 光滑），由\eqref{align:ch09-19}推出 $\Psi^{(t_0;p)}(t)=\Psi(t;t_0;p)$ 是具有初始条件 $\Psi^{(t_0;p)}(t_0)=p$ 的 $V$ 的积分曲线。

为证明唯一性，假设 $t_0\in J$，$\gamma\colon J_0\to M$ 是定义在包含 $t_0$ 的某开区间 $J_0\subseteq J$ 上、满足 $\gamma(t_0)=p$ 的任意 $V$ 的积分曲线。定义光滑曲线 $\widetilde{\gamma}\colon J_0\to J\times M$ 为 $\widetilde{\gamma}(t)=(t;\gamma(t))$。易见 $\widetilde{\gamma}$ 是具有初始条件 $\widetilde{\gamma}(t_0)=(t_0;p)$ 的 $\widetilde{V}$ 的积分曲线。由 $\widetilde{V}$ 积分曲线的唯一性和极大性，在其整个定义域上必有 $\widetilde{\gamma}(t)=\widetilde{\theta}(t-t_0;(t_0;p))$，这推出 $\gamma$ 的定义域包含于 $\Psi^{(t_0;p)}$ 的定义域，且在该定义域上 $\gamma=\Psi^{(t_0;p)}$。由此 $\Psi^{(t_0;p)}$ 是在 $t=t_0$ 经过 $p$ 的 $V$ 的唯一极大积分曲线。这完成了 (a) 的证明。

为证明 (b)，假设 $t_1\in E(t_0;p)$，$q=\Psi^{(t_0;p)}(t_1)$。则 $\Psi^{(t_1;q)}$ 和 $\Psi^{(t_0;p)}$ 都是在 $t=t_1$ 经过 $q$ 的 $V$ 的积分曲线，故由唯一性和极大性它们必有相同定义域且在其上相等。

接下来证明 (d)。假设 $p\in M_{t_1;t_0}$，$\Psi_{t_1;t_0}(p)\in M_{t_2;t_1}$，令 $q=\Psi_{t_1;t_0}(p)=\Psi^{(t_0;p)}(t_1)$。则 (b) 推出 $\Psi^{(t_1;q)}(t_2)=\Psi^{(t_0;p)}(t_2)$。展开定义即得\eqref{align:ch09-18}。

最后证明 (c)。假设 $(t_1;t_0)\in J\times J$。我们已注意到 $M_{t_1;t_0}$ 在 $M$ 中开。为证明 $\Psi_{t_1;t_0}(M_{t_1;t_0})\subseteq M_{t_0;t_1}$，令 $p$ 为 $M_{t_1;t_0}$ 的点，令 $q=\Psi_{t_1;t_0}(p)$。(b) 推出 $E(t_0;p)=E(t_1;q)$，故 $t_0\in E(t_0;p)=E(t_1;q)$。这等价于 $(t_0;t_1;q)\in E$，从而如所断言 $q\in M_{t_0;t_1}$。为看到 $\Psi_{t_1;t_0}\colon M_{t_1;t_0}\to M_{t_0;t_1}$ 是微分同胚，只需注意上述同样的论证推出 $\Psi_{t_0;t_1}(M_{t_0;t_1})\subseteq M_{t_1;t_0}$，然后 (d) 推出对所有 $p\in M_{t_0;t_1}$，$\Psi_{t_1;t_0}\circ\Psi_{t_0;t_1}(p)=\Psi_{t_1;t_1}(p)=p$，类似地对 $q\in M_{t_1;t_0}$，$\Psi_{t_0;t_1}\circ\Psi_{t_1;t_0}(q)=q$。
\end{proof}

\begin{exercise}\label{exercise:ch09-49}
设 $M$ 是光滑流形。假设 $X$ 是 $M$ 上的（不含时）光滑向量场，$\theta\colon\mathscr{D}\to M$ 是其流。设 $V$ 是由 $V(t;p)=X_p$ 定义的含时向量场。证明 $V$ 的含时流由 $\Psi(t;t_0;p)=\theta(t-t_0;p)$ 给出，定义域 $E=\{(t;t_0;p):(t-t_0;p)\in\mathscr{D}\}$。
\end{exercise}

\begin{example}\label{example:ch09-50}
在 $\mathbb{R}^n$ 上定义含时向量场 $V$ 为
\begin{align*}
V(t;x)=\frac1t\,x^i\frac{\partial}{\partial x^i}\bigg|_x,\qquad (t;x)\in(0;\infty)\times\mathbb{R}^n.
\end{align*}
任取 $t_0\in(0;\infty)$，$x_0\in\mathbb{R}^n$，令 $\gamma(t)=(x^1(t),\dots,x^n(t))$ 表示具有初始条件 $\gamma(t_0)=x_0$ 的 $V$ 的积分曲线。则 $\gamma$ 的分量满足如下非自治微分方程组：
\begin{align*}
\dot{x}^i(t)=\frac1t x^i(t),\qquad x^i(t_0)=x^i_0.
\end{align*}
容易验证此方程组的极大解为 $x^i(t)=tx^i_0/t_0$，对所有 $t>0$ 有定义。因此，$V$ 的含时流为 $\Psi(t;t_0;x)=tx/t_0$，对 $(t;t_0;x)\in(0;\infty)\times(0;\infty)\times\mathbb{R}^n$。
\end{example}

\end{document}